% Source: Corsin Nick/Class Notes/Week 6.pdf, p. 1
\chapter{Week 6 --- Convexity}

\exercisesheet{6}

\textit{Statements quoted verbatim from \texttt{exercises/Ex6\_Analysis2\_eng.pdf} (assigned 20
March 2026, due 30 March 2026). Attribution: Prof. Joaquim Serra, D-MATH, ETH Z\"urich.}

\begin{ainote}[Corsin's recommendations]
Colour key, as given on the page: blue \textbf{= important}, orange \textbf{= semi-important},
red \textbf{= optional}. 6.1 and 6.4 carry no colour marker on the page.

\begin{center}
\begin{tabular}{lll}
\textbf{Problem} & \textbf{Priority} & \textbf{Corsin's note} \\
\hline
6.1 & unmarked & --- \\
6.2 & important & ``This is what we covered on Friday. Find a Lemma in your linear \\
    & & algebra notes which relates the determinant and trace to the eigenvalues.'' \\
6.3 & semi-important & ``By using $\alpha+\beta+\gamma=\pi$ for the angles, you can express \\
    & & this as an optimization problem in two variables.'' \\
6.4 & unmarked & --- \\
6.5 & optional & ``See `data analysis'.'' \\
6.6 & optional & ``Do two or three, not all of them.'' \\
6.7 & important & --- \\
6.8 & important & --- \\
6.9 & important & --- \\
6.10 & important & \multirow{3}{*}{\shortstack[l]{``These are good exercises, but how many\\ optimization problems do you really need?''}} \\
6.11 & important & \\
6.12 & important & \\
\end{tabular}
\end{center}
\end{ainote}

\begin{ainote}
Only the exercises tagged \textbf{important} above are reproduced below (their wording only, no
solutions yet): 6.2, 6.7--6.9, 6.10--6.12. The full sheet (6.1--6.12) has not yet been
transcribed into \texttt{content/exercise-sheets/week-06.tex}.
\end{ainote}

\begin{exercise}[6.2 --- The signature of a $2\times 2$ matrix \textnormal{(important)}]
\label{ex:6.2}
Despite the definition, it is not necessary to compute the eigenvalues of a matrix to find its
signature. Prove that for a symmetric $2\times 2$ matrix $M$ we have the following simple rule to
determine the signature in terms of $\det M$ and $\Tr M$:
\begin{itemize}
  \item If $\det M > 0$, $\Tr M > 0$ then $M$ is positive definite,
  \item If $\det M > 0$, $\Tr M < 0$ then $M$ is negative definite,
  \item If $\det M < 0$ then $M$ is indefinite,
  \item If $\det M = 0$, then $M$ is degenerate.
\end{itemize}
\end{exercise}

\begin{exercise}[6.7 --- Multiple choice \textnormal{(important)}]
\label{ex:6.7}
Among the following statements about convex functions mark those (and only those) which are
always true.
\begin{enumerate}[label=\textbf{(\alph*)}]
  \item If $f \in C^1(U)$ is convex in some open convex set $U \subset \mathbb{R}^n$ and $f$ has a
    local maximum at $z \in U$, then $\nabla f \equiv 0$ in $U$.
  \item If $f \in C^1(U)$ is convex in some open convex set $U \subset \mathbb{R}^n$ and $f$ has a
    global maximum at $z \in U$, then $\nabla f \equiv 0$ in $U$.
  \item Assume $f_n \in C^2(\mathbb{R})$ is a sequence of convex functions that converge
    pointwise to some $f : \mathbb{R} \to \mathbb{R}$. Is $f$ necessarily convex?
  \item There exists a convex function $f \in C^\infty(\mathbb{R}^2)$ such that
    $$f(x) = 1 - 2x_1 + x_2^3 + O(|x|^4) \quad \text{as } |x| \to 0.$$
  \item There exists a convex function $f \in C^\infty(\mathbb{R}^2)$ such that
    $$f(x) = 1 - 2x_1 + x_2^4 + O(|x|^4) \quad \text{as } |x| \to 0.$$
  \item A convex set is not necessarily connected.
\end{enumerate}
\end{exercise}

\begin{exercise}[6.8 --- Multiple choice \textnormal{(important)}]
\label{ex:6.8}
The Hessian matrix of $f \in C^2(\mathbb{R}^n)$ is positive semidefinite at a critical point $x_0$
of $f$, i.e.,
$$\langle v, \Hess f(x_0) v\rangle \geq 0 \quad \text{for all } v \in \mathbb{R}^n.$$
Which of the following statements necessarily hold? (There may be more than one).
\begin{enumerate}[label=\textbf{(\alph*)}]
  \item $x_0$ is a strict local minimum of $f$.
  \item $x_0$ is a local minimum of $f$.
  \item $x_0$ is not a local maximum of $f$.
  \item None of the above statements.
\end{enumerate}
\end{exercise}

\begin{exercise}[6.9 --- Minimization \textnormal{(important)}]
\label{ex:6.9}
The function $f : \mathbb{R}^2 \to \mathbb{R}$ is given by $f(x,y) = 2x^2+y^2-x$. Determine the
extrema of $f$ on\dots
\begin{enumerate}[label=\textbf{(\alph*)}]
  \item \dots the unit circle $S^1 = \{(x,y) \in \mathbb{R}^2 \mid x^2+y^2 = 1\}$;
  \item \dots the closed unit disk $D = \{(x,y) \in \mathbb{R}^2 \mid x^2+y^2 \leq 1\}$.
\end{enumerate}
\end{exercise}

\begin{exercise}[6.10 --- Lagrange multipliers \textnormal{(important)}]
\label{ex:6.10}
Consider the function $f(x,y,z) = 3x-y+2z$ and the set
$$M = \{(x,y,z) \in \mathbb{R}^3 \mid x^2+y^2+z^2 = 1,\ x+y = 0\}.$$
Determine the extrema of $f$ on $M$ and their nature.
\end{exercise}

\begin{exercise}[6.11 --- Closest point on a hyperboloid \textnormal{(important)}]
\label{ex:6.11}
Find the points on the submanifold
$$M = \{(x,y,z) \in \mathbb{R}^3 : z = x^2-y^2\}$$
that are closest to the point $(0,0,1)$.
\end{exercise}

\begin{exercise}[6.12 --- Critical points \textnormal{(important)}]
\label{ex:6.12}
Let $f : \mathbb{R}^2 \to \mathbb{R}$ be the function $f(x,y) = (ax^2+by^2)e^{-x^2-y^2}$ with real
parameters $a, b \in \mathbb{R}$. Find all critical points and determine their nature with the
Hessian test, depending on $a, b$.
\end{exercise}

\section{Solutions}
\label{sec:week06_solutions}

\begin{ainote}
Corsin does not work out any of the official-sheet exercises in his own notes for this week (his
Week 6 PDF covers only convexity and the inverse/implicit function theorems, with no page solving
6.2 or 6.7--6.12) so all solutions below are original. Each has been checked against
\texttt{exercises/Sol6\_Analysis2\_eng.pdf}; no divergence from the official solutions was found.
\end{ainote}

\begin{exercisesolution}[Proof of \cref{ex:6.2}]
% Generator: Claude Sonnet 5 (Medium)
By the spectral theorem, since $M$ is symmetric there exists an orthogonal matrix $O$
(so $OO\transp = \id$, hence $\det O = \pm 1$) and $\lambda_1, \lambda_2 \in \mathbb{R}$ such that
$$OMO\transp = \begin{pmatrix} \lambda_1 & 0 \\ 0 & \lambda_2 \end{pmatrix}.$$
Using $\det(AB) = \det(A)\det(B)$ and $\Tr(AB) = \Tr(BA)$,
$$\det M = \det(OMO\transp) = \lambda_1\lambda_2, \qquad \Tr M = \Tr(MO\transp O) = \Tr(OMO\transp) = \lambda_1 + \lambda_2.$$
So $\det M$ and $\Tr M$ are exactly the product and sum of the eigenvalues, and the four cases
follow immediately:
\begin{itemize}
  \item $\det M > 0$: $\lambda_1, \lambda_2$ have the same sign. If additionally $\Tr M > 0$,
    both are positive, so $M$ is \textbf{positive definite}; if $\Tr M < 0$, both are negative,
    so $M$ is \textbf{negative definite}.
  \item $\det M < 0$: $\lambda_1, \lambda_2$ have opposite signs, so $M$ is \textbf{indefinite}.
  \item $\det M = 0$: one of $\lambda_1, \lambda_2$ is zero, so $M$ is \textbf{degenerate}.
\end{itemize}
\end{exercisesolution}

\begin{exercisesolution}[Solution to \cref{ex:6.7}]
% Generator: Claude Sonnet 5 (Medium)
\begin{enumerate}[label=\textbf{(\alph*)}]
  \item \textbf{False.} Let $f(x) := \max\{x_1, \tfrac12\}^4$ on $U := B_1 \subset \mathbb{R}^n$.
    Being $(\max\{x_1,\tfrac12\})^4$ with $\max\{x_1,\tfrac12\} \geq \tfrac12 > 0$, and $t \mapsto
    t^4$ convex and increasing on $(0,\infty)$ composed with the convex function
    $x \mapsto \max\{x_1,\tfrac12\}$, $f$ is convex. For all $x \in U$ with $x_1 \leq \tfrac12$,
    $f(x) = (\tfrac12)^4$ identically, so $f$ has a (non-strict) local maximum at $x=0$, yet
    $\nabla f$ does not vanish once $x_1 > \tfrac12$.

    \begin{ainote}
    Worth noting: $x=0$ here is simultaneously a (non-strict) local \emph{minimum} too, since $f$
    is locally constant on the whole region $\{x_1 \leq \tfrac12\} \cap U$. This is exactly why
    the counterexample doesn't contradict \textbf{(b)}: $x=0$ is not a \emph{global} maximum of
    $f$ on $U$ (as $f$ grows without bound as $x_1 \to 1$), so \textbf{(b)}'s stronger hypothesis
    simply isn't triggered.
    \end{ainote}
  \item \textbf{True.} Since $z$ is a local maximum, $\nabla f(z) = 0$
    (\cref{prop:extremum_implies_critical}). Convexity gives $f(x) \geq f(z) + \nabla f(z)\cdot(x-z)
    = f(z)$ for all $x \in U$; but $z$ being a global maximum also gives $f(x) \leq f(z)$ for all
    $x \in U$. Hence $f \equiv f(z)$ on $U$, so $\nabla f \equiv 0$ in $U$.
  \item \textbf{True.} Fix $x, y \in U$ and $t \in [0,1]$. Convexity of each $f_n$ gives
    $f_n(tx+(1-t)y) \leq tf_n(x) + (1-t)f_n(y)$ for every $n$; letting $n \to \infty$ (with
    $x,y,t$ fixed) and using pointwise convergence yields the same inequality for $f$. Since
    $x,y,t$ were arbitrary, $f$ is convex.
  \item \textbf{False.} A convex function lies above its tangent plane, so from the Taylor data at
    $0$, any such $f$ would satisfy $f(x) \geq 1-2x_1$ for all $x \in \mathbb{R}^n$. Combined with
    $f(x) = 1-2x_1+x_2^3+O(|x|^4)$, this forces $x_2^3 \geq -M|x|^4$ near $0$ for some $M>0$,
    which fails along $x = (0,-r,0,\dots,0)$ as $r \downarrow 0$ (left side $\sim -r^3$, right
    side $\sim -Mr^4$, and $-r^3 < -Mr^4$ for small $r$). Contradiction, so no such convex $f$
    exists.
  \item \textbf{True.} Simply take $f(x) := 1-2x_1+x_2^4$ itself (with the $O(|x|^4)$ remainder
    identically zero): it is convex, being a sum of the affine (hence convex) term $1-2x_1$ and
    the convex term $x_2^4$.
  \item \textbf{False.} A convex set is always path-connected, since for any two points the
    straight segment joining them lies inside the set by definition of convexity; path-connected
    sets are connected.
\end{enumerate}
\end{exercisesolution}

\begin{exercisesolution}[Solution to \cref{ex:6.8}]
% Generator: Claude Sonnet 5 (Medium)
None of \textbf{(a)}--\textbf{(c)} hold in general:
\begin{itemize}
  \item \textbf{(a) False:} $f \equiv 0$ has $\Hess f(x_0) = 0$ positive semidefinite at every
    $x_0$, but $x_0$ is not a \emph{strict} local minimum (every nearby value equals $f(x_0)$).
  \item \textbf{(b) False:} $f(x) = x_1^3$ at $x_0 = 0$ has $\Hess f(0) = 0$ (positive
    semidefinite), but $f$ takes negative values arbitrarily close to $0$, so $x_0$ is not a
    local minimum at all.
  \item \textbf{(c) False:} $f(x) = -x_1^4$ at $x_0 = 0$ has $\Hess f(0) = 0$ (positive
    semidefinite, vacuously, since it is the zero matrix), yet $x_0$ \emph{is} a (strict) local
    maximum of $f$.
\end{itemize}
Since \textbf{(a)}--\textbf{(c)} all fail, \textbf{(d) is the correct answer}: none of the above
statements necessarily hold.
\end{exercisesolution}

\begin{exercisesolution}[Solution to \cref{ex:6.9}]
% Generator: Claude Sonnet 5 (Medium)
\begin{enumerate}[label=\textbf{(\alph*)}]
  \item With $g(x,y) := x^2+y^2-1$, the Lagrangian is $L(x,y,\lambda) = f(x,y) - \lambda g(x,y)$,
    giving the system
    $$0 = \partial_x L = 2x(2-\lambda) - 1, \qquad 0 = \partial_y L = 2y(2-2\lambda), \qquad 0 = \partial_\lambda L = -(x^2+y^2-1).$$
    From the middle equation, either $y = 0$ or $\lambda = 1$.

    \textbf{If $y = 0$:} then $x^2 = 1$, so $x = \pm 1$, giving the points $(1,0)$ and $(-1,0)$
    with $f(1,0) = 1$ and $f(-1,0) = 3$.

    \textbf{If $\lambda = 1$:} the first equation gives $2x(1) = 1$, i.e.\ $x = \tfrac12$, and the
    constraint gives $y = \pm\tfrac{\sqrt3}{2}$, with $f\big(\tfrac12,\pm\tfrac{\sqrt3}{2}\big) =
    \tfrac12 + \tfrac34 - \tfrac12 = \tfrac34$.

    Comparing all four values, $f$ attains a \textbf{global maximum} at $(-1,0)$ (value $3$) and
    \textbf{global minima} at $\big(\tfrac12,\pm\tfrac{\sqrt3}{2}\big)$ (value $\tfrac34$); $(1,0)$
    is a local (non-global) extremum with value $1$.
  \end{enumerate}
\end{exercisesolution}

\begin{ainote}
The official solution lists $(1,0)$ alongside $(-1,0)$ as "local extrema" from the Lagrange
system, then only mentions it again implicitly by omission from the final global-extrema summary
--- worth flagging explicitly: \textbf{Lagrange's condition finds all \emph{critical} points of
$f|_{S^1}$, not automatically the global extrema}. Here $(1,0)$ is a genuine local maximum along
the circle (value $1$) that is simply not the largest one; only comparing all candidate values
at the end distinguishes it from the global maximum at $(-1,0)$ (value $3$).
\end{ainote}

\begin{exercisesolution}[Solution to \cref{ex:6.9}, continued]
% Generator: Claude Sonnet 5 (Medium)
\begin{enumerate}[label=\textbf{(\alph*)}, start=2]
  \item Having handled the boundary $S^1 = \partial D$ in \textbf{(a)}, we look for interior
    critical points: $Df(x,y) = (4x-1, 2y) = 0$ gives $(x,y) = \big(\tfrac14, 0\big) \in
    \operatorname{int}(D)$. The Hessian $\Hess f = \begin{pmatrix} 4 & 0 \\ 0 & 2\end{pmatrix}$ is
    positive definite, so this is a local minimum, with value $f\big(\tfrac14,0\big) =
    \tfrac18 - \tfrac14 = -\tfrac18$. Since $-\tfrac18$ is smaller than every boundary value found
    in \textbf{(a)}, $f$ on $D$ attains a \textbf{global minimum} at $\big(\tfrac14,0\big)$ and a
    \textbf{global maximum} at $(-1,0)$.
\end{enumerate}
\end{exercisesolution}

\begin{exercisesolution}[Solution to \cref{ex:6.10}]
% Generator: Claude Sonnet 5 (Medium)
With $g_1(x,y,z) := x^2+y^2+z^2-1$ and $g_2(x,y,z) := x+y$, the Lagrangian
$L = f - \lambda_1 g_1 - \lambda_2 g_2$ gives the system
$$3 - 2\lambda_1 x - \lambda_2 = 0, \qquad -1-2\lambda_1 y - \lambda_2 = 0, \qquad 2 - 2\lambda_1 z = 0, \qquad x+y=0, \qquad x^2+y^2+z^2=1.$$
From $x+y=0$, $y=-x$. Subtracting the second equation from the first eliminates $\lambda_2$:
$$4 - 2\lambda_1 x - 2\lambda_1(-x) \cdot(-1)= 4-2\lambda_1x-2\lambda_1x=4-4\lambda_1 x = 0 \implies \lambda_1 x = 1.$$
From $2-2\lambda_1 z=0$, $\lambda_1 = 1/z$ (so $z \neq 0$), hence $x = z$ (from $\lambda_1 x = 1 =
\lambda_1 z$). With $y=-x=-z$, the constraint becomes $3x^2 = 1$, i.e.\ $x = \pm\tfrac{1}{\sqrt3}$,
giving the two points $\tfrac{1}{\sqrt3}(1,-1,1)$ and $-\tfrac{1}{\sqrt3}(1,-1,1)$. Evaluating,
$$f\Big(\tfrac{1}{\sqrt3}(1,-1,1)\Big) = \frac{3+1+2}{\sqrt3} = 2\sqrt3, \qquad f\Big({-\tfrac{1}{\sqrt3}}(1,-1,1)\Big) = -2\sqrt3.$$
Since $M$ is compact (closed and bounded in $\mathbb{R}^3$) and $f$ is continuous, these are the
only Lagrange candidates, so $f$ attains its \textbf{maximum} $2\sqrt3$ at $\tfrac{1}{\sqrt3}(1,-1,1)$
and its \textbf{minimum} $-2\sqrt3$ at $-\tfrac{1}{\sqrt3}(1,-1,1)$.
\end{exercisesolution}

\begin{exercisesolution}[Solution to \cref{ex:6.11}]
% Generator: Claude Sonnet 5 (Medium)
Minimizing the distance to $(0,0,1)$ is the same as minimizing
$f(x,y,z) := x^2+y^2+(z-1)^2$ subject to $g(x,y,z) := z-x^2+y^2 = 0$. Since $f$ is a distance
squared to a fixed point, its minimum over the (closed) constraint set is attained. Lagrange's
condition $\nabla f = \lambda \nabla g$ reads
$$2x = -2\lambda x, \qquad 2y = 2\lambda y, \qquad 2(z-1) = \lambda,$$
i.e.\ $x(1+\lambda) = 0$ and $y(1-\lambda) = 0$.

\textbf{Case $x=y=0$:} the constraint gives $z=0$, so $(0,0,0)$ with $d^2 = 1$.

\textbf{Case $x \neq 0$ (so $\lambda=-1$):} then $2(z-1)=-1 \implies z = \tfrac12$, and $y(1-\lambda)=2y=0
\implies y=0$. The constraint $z=x^2-y^2$ gives $x^2=\tfrac12$, i.e.\ $x=\pm\tfrac{1}{\sqrt2}$,
giving $\big(\pm\tfrac{1}{\sqrt2},0,\tfrac12\big)$ with $d^2 = \tfrac12+0+\tfrac14 = \tfrac34$.

\textbf{Case $y \neq 0$ (so $\lambda=1$):} then $2(z-1)=1 \implies z=\tfrac32$, and $x(1+\lambda)=2x=0
\implies x=0$; the constraint becomes $-y^2 = \tfrac32$, which has no real solution.

Comparing $d^2=1$ against $d^2=\tfrac34$, the \textbf{closest points} are
$\big(\pm\tfrac{1}{\sqrt2},0,\tfrac12\big)$, at minimal distance $\tfrac{\sqrt3}{2}$.
\end{exercisesolution}

\begin{exercisesolution}[Solution to \cref{ex:6.12}]
% Generator: Claude Sonnet 5 (Medium)
Writing $r^2 := x^2+y^2$, the gradient is
$$\nabla f(x,y) = \big(2x(a-ax^2-by^2),\ 2y(b-ax^2-by^2)\big)e^{-r^2}.$$

\textbf{Case $a=b=0$:} $f \equiv 0$, so every point of $\mathbb{R}^2$ is critical, and (being the
constant function) each is simultaneously a global maximum and minimum.

\textbf{Case $a=0, b\neq 0$ (and symmetrically $a \neq 0, b=0$):} the gradient equations reduce to
$-2xby^2=0$ and $2by(1-y^2)=0$, whose common solutions are the entire line $\{(x,0): x\in
\mathbb{R}\}$ together with the two points $(0,\pm1)$. On the critical line the Hessian is
singular (degenerate, as expected since a whole line is critical); at $(0,\pm1)$,
$\Hess f(0,\pm1) = \operatorname{diag}(-2b,-4b)\,e^{-1}$, giving \textbf{minima} if $b<0$ and
\textbf{maxima} if $b>0$. The case $a\neq0,b=0$ is symmetric, with critical line $\{(0,y)\}$ and
extra points $(\pm1,0)$: minima if $a<0$, maxima if $a>0$.

\textbf{Case $a,b \neq 0$, $a \neq b$:} setting $x=0$ or $y=0$ in the gradient equations gives the
five critical points $(0,0)$, $(0,\pm1)$, $(\pm1,0)$ (checking $x,y\neq0$ simultaneously forces
$a=b$, excluded here). The Hessians are
$$\Hess f(0,0) = \begin{pmatrix}2a&0\\0&2b\end{pmatrix}, \quad \Hess f(0,\pm1) = \begin{pmatrix}2(a-b)&0\\0&-4b\end{pmatrix}e^{-1}, \quad \Hess f(\pm1,0) = \begin{pmatrix}-4a&0\\0&2(b-a)\end{pmatrix}e^{-1}.$$
Reading off signs for each ordering of $a,b$ against $0$:
\begin{itemize}
  \item $a>b>0$: min at $(0,0)$, saddle at $(0,\pm1)$, max at $(\pm1,0)$.
  \item $b>a>0$: min at $(0,0)$, max at $(0,\pm1)$, saddle at $(\pm1,0)$.
  \item $a>0>b$: saddle at $(0,0)$, min at $(0,\pm1)$, max at $(\pm1,0)$.
  \item $b>0>a$: saddle at $(0,0)$, max at $(0,\pm1)$, min at $(\pm1,0)$.
  \item $0>a>b$: max at $(0,0)$, min at $(0,\pm1)$, saddle at $(\pm1,0)$.
  \item $0>b>a$: max at $(0,0)$, saddle at $(0,\pm1)$, min at $(\pm1,0)$.
\end{itemize}

\textbf{Case $a=b\neq0$:} beyond $(0,0)$, $(0,\pm1)$, $(\pm1,0)$, setting $x,y\neq0$ now gives
$a-ax^2-ay^2=0 \iff x^2+y^2=1$, so the entire unit circle is critical. There,
$\Hess f(x,y) = \operatorname{diag}(-4ax^2,-4ay^2)$ off-diagonal $-4axy$, which has determinant
$0$ (singular), consistent with a whole curve of critical points.
\end{exercisesolution}
